\section{第\ref{sec:neurontypes}章　神经元类型}

本章的数字在有直接计数的地方依据人脑细胞的直接计数；“一个神经元，一种递质”的规则依据细胞图谱以及发现了例外的实验；\nt{DOPA}的作用依据脉冲记录；其余广播通道的作用则依据假说。

{\footnotesize
\setlength{\tabcolsep}{3pt}\renewcommand{\arraystretch}{1.15}
\begin{longtable}{>{\raggedright}p{0.5cm}>{\raggedright}p{4.0cm}>{\raggedright}p{5.6cm}>{\raggedright}p{2.3cm}>{\raggedright\arraybackslash}p{1.7cm}}
\toprule
编号 & 命题 & 实验与测量内容 & 来源 & 状态 \\
\midrule\endfirsthead
\toprule
编号 & 命题 & 实验与测量内容 & 来源 & 状态 \\
\midrule\endhead
1 & 人脑约有 $8.6\cdot10^{10}$ 个神经元，\nt{GLUT}神经元几乎全是小脑的小神经元（表~\ref{tab:types}） & 各向同性分离法：把组织溶解成均匀的细胞核悬液，计数带神经元标记物NeuN的细胞核。成年男性有 $(86.1\pm8.1)\times10^9$ 个神经元；其中只有 $19\,\%$ 位于大脑皮层，大部分在小脑 & \cite{Azevedo2009} & 已测量 \\
2 & 几乎每个神经元只有一种主要递质（命题~\ref{prop:dale}），但有例外 & 小鼠脑转录组图谱：约 $4\cdot10^6$ 个细胞，$5322$ 个簇；按合成与包装递质的基因为各簇指定递质。例外有直接测量：脑片中的\nt{DOPA}神经元经同一种囊泡转运体同时释放GABA，抑制纹状体神经元；部分\nt{DOPA}轴突的谷氨酸和多巴胺从同一根导线的不同末梢释放（电子显微镜和光遗传学） & \cite{Strata1999,Yao2023,Tritsch2012,Zhang2015} & 已测量 \\
3 & 整个权重矩阵列同号~\eqref{eq:dalesign} & 本章由命题~\ref{prop:dale} 以及“谷氨酸受体几乎都是兴奋性的、GABA受体几乎都是抑制性的”推出。“同一神经元的所有接收者得到同号权重”作为一般规则没有直接检验；反例是\nt{DOPA}神经元共同释放GABA（第~2行）和带有不同符号受体的接收者（第~11行） & 本章推导 & 理论 \\
4 & 皮层神经元中四分之三到五分之四是\nt{GLUT}，五分之一到四分之一是\nt{GABA} & 在猴皮层 $10$ 个区中，对贯穿皮层全层、宽 $50\um$ 的柱条做GABA免疫染色：$8$ 个区中GABA神经元约占全部神经元的 $25\,\%$，初级视皮层中不到 $20\,\%$ & \cite{Hendry1987} & 已测量 \\
5 & 一个\nt{GLUT}神经元约有 $10^4$ 个输出 & 尸检体视学：年轻男性新皮层中有 $1.64\cdot10^{14}$ 个突触（电子显微镜，体视框法）和约 $2.3\cdot10^{10}$ 个神经元。两者之比为每个神经元 $\sim7\cdot10^3$ 个突触；平均而言输入数等于输出数 & \cite{Tang2001synapses,Pakkenberg1997} & 已测量 \\
6 & \nt{DOPA}神经元的输出分支成 $10^5$--$10^6$ 个末梢；这是根据靶区覆盖范围的估计，而非计数 & 用病毒标记大鼠单个\nt{DOPA}神经元的膜，完整重建轴突：一个神经元覆盖纹状体体积的 $0.45$ 到 $5.7\,\%$（平均 $2.7\,\%$）。测量的是覆盖范围，而非末梢数：末梢数由覆盖范围按量级推出，没有直接计数。\nt{SERO}和\nt{NORA}的末梢数是粗略估计 & \cite{Matsuda2009} & 已测量 \\
7 & 广播神经元很少：\nt{DOPA} $\sim5\cdot10^5$（两群中的主要一群，下限），\nt{SERO} $\sim3\cdot10^5$（两次局部计数之和），\nt{HIST} $\sim6\cdot10^4$（表~\ref{tab:types}，图~\ref{fig:typepie}） & 人体体视学计数：黑质中 $550\,000$ 个含色素神经元（不含腹侧被盖区）；中缝背核 $235\,000$ 个神经元（该核全部神经元），脑桥被盖另有 $125\,000$ 个血清素神经元；下丘脑约 $64\,000$ 个组胺神经元。\nt{NORA}、\nt{GLYC}、\nt{ACET}和\nt{ADRE}没有直接计数：表中是粗略的量级估计 & \cite{Pakkenberg1991,Baker1990,Baker1991,Airaksinen1991} & 已测量 \\
8 & 缝隙中的谷氨酸猛增到约一毫摩尔，在 $\sim1\ms$ 内下降 & 根据突触传递过程中快速解离的竞争性阻断剂从NMDA受体上被置换的情况（海马神经元培养物）：峰值 $1.1$\,mM，衰减时间常数 $1.2\ms$ & \cite{Clements1992} & 已测量 \\
9 & 工作中的皮层里兴奋性电流与抑制性电流几乎平衡，神经元停在阈值附近 & 理论：强连接网络会自行进入平衡状态。实验：雪貂前额叶皮层在体记录中，在“上”状态期间突触电流的反转电位在数百毫秒内保持在 $-37$\,mV左右，而电导变化达 $\sim21$\,nS：兴奋与抑制同升同降 & \cite{vanVreeswijk1996,Haider2006} & 已测量 \\
10 & \nt{DOPA}神经元报告奖励预测误差~\eqref{eq:rpe}（图~\ref{fig:rpe}） & 猴\nt{DOPA}神经元的脉冲：对意外的果汁发出簇放电；学习后对提示信号发出簇放电，对被预测的果汁不响应；许诺的果汁没给时放电频率下陷。在定量实验中，频率与当前奖励和过去奖励加权平均之差成正比，但只对正误差如此。方程~\eqref{eq:rpe} 本身是时间差分算法 & \cite{Schultz1997,Bayer2005,SuttonBarto2018} & 已测量 \\
11 & 符号和速度由受体决定：离子型不到一毫秒，代谢型慢得多（命题~\ref{prop:receptor}） & 向海马神经元的膜片施加短促的谷氨酸脉冲：经离子型受体的电流在 $0.2$--$0.6\ms$ 内上升（从 $20$ 到 $80\,\%$），在 $2$--$3\ms$ 内衰减。锥体神经元的慢抑制电位可被代谢型GABA受体的选择性阻断剂消除。多巴胺受体有两个作用相反的家族；在纹状体中，带D1受体的神经元需要多巴胺来增强突触，带D2受体的神经元需要多巴胺来减弱突触 & \cite{Colquhoun1992,DutarNicoll1988,Beaulieu2011,Shen2008} & 已测量 \\
12 & 广播通道不是一根导线，而是由少数线路组成的线束（\ref{sec:buses}~节） & 用病毒遗传学方法标记小鼠中缝背核的血清素神经元：投射到杏仁核和投射到额叶皮层的神经元位于核内不同位置，分支到互补的区域，对厌恶刺激的反应相反。综述：蓝斑（\nt{NORA}）的神经元亚群编码不同的过程 & \cite{Ren2018,Poe2020} & 已测量 \\
13 & \nt{NORA}设定增益 $g$ & 自适应增益假说；儿茶酚胺提高传递特性陡度的网络模型。“$g$ 随去甲肾上腺素增大”在整个网络中没有直接测量；已测得瞳孔直径跟随猴蓝斑神经元的脉冲变化 & \cite{AstonJones2005,ServanSchreiber1990,Joshi2016} & 假说 \\
14 & \nt{ACET}切换“写入/读取”并设定学习速率 & 假说。依据：在大鼠嗅皮层脑片中，乙酰胆碱激动剂（$100\,\mu$M）使内部“回忆”突触减弱超过 $60\,\%$，而使输入突触减弱不到 $15\,\%$ & \cite{Hasselmo2006,HasselmoBower1992} & 假说 \\
15 & \nt{SERO}设定折扣因子 $\gamma$；血清素神经元放电平稳，每秒几个脉冲，在睡眠的某一阶段几乎沉默 & “血清素 $=\gamma$”假说。最有力的论据：光遗传学激活小鼠中缝背核的血清素神经元，会减少在等待延迟奖励时过早放弃的次数。脉冲频率以及快速眼动睡眠中的沉默是测量过的（记录综述） & \cite{Doya2002,Miyazaki2014,JacobsAzmitia1992} & 假说 \\
\bottomrule
\end{longtable}}
